\begin{textsample}{POS Dim 4 – 2015\_c6 – Score 14.00 – 2015\_c6/enade\_2015\_admin\_8d6df9fa83.txt}  \label{ex:f4_pos_023}
CAPÍTULO 3 ANÁLISE TÉCNICA DA PROVA Este capítulo tem por objetivo apresentar o desempenho dos estudantes concluintes de Administração no Enade/2015.
Para isso, foram calculadas as estatísticas básicas da prova em seu todo ( seção 3.1.1 ), bem como as estatísticas relacionadas aos Componentes de Formação Geral ( seção 3.1.2 ) e de Conhecimento Específico da Área ( seção 3.1.3 ).
Dadas as suas características, foram analisadas, em separado, as questões objetivas ( seção 3.2 ) e as questões discursivas ( seção 3.3 ).
Para as questões objetivas, foram disponibilizados os índices de facilidade e de discriminação ponto-bisserial, também em separado, para os Componentes de Formação Geral ( seção 3.2.1 ) e de Conhecimento Específico ( seção 3.2.2 ).
De cada componente, uma das questões foi escolhida para exemplificar a análise gráfica, relacionando as alternativas escolhidas pelos estudantes ( inclusive o gabarito ) com o número de acertos no componente.
O Anexo I apresenta a íntegra da análise gráfica para todas as questões objetivas.
Para cada uma das questões discursivas, os conteúdos dos tipos mais comuns de respostas dos estudantes são apresentados e comparados com o padrão de respostas esperado ( ver Anexo VIII com o padrão de respostas ).
Tomando como base as duas questões discursivas do Componente de Formação Geral, a seção 3.3.1.6 apresenta comentários sobre a correção das respostas com respeito à Língua Portuguesa.
Nas Tabelas 3.1 e 3.3 são apresentados o tamanho da população inscrita, ausente, presente e percentual de faltosos, respectivamente, por Grande Região e por Categoria Administrativa e Organização Acadêmica da IES.
As demais tabelas apresentam as seguintes estatísticas das notas¹⁵ : média do desempenho na prova, erro padrão da média, desvio padrão, nota mínima, mediana e nota máxima para cada um de seus componentes.
Tais estatísticas contemplam o total de estudantes concluintes da Área de Administração inscritos e presentes à prova do Enade/2015, tendo em vista agregações ou por Grandes Regiões e o país como um todo, ou por Categoria Administrativa e Organização Acadêmica da IES.
Em relação aos gráficos de distribuição de notas, o intervalo considerado foi de 10 unidades, aberto à esquerda e fechado à direita, com exceção do primeiro intervalo, [ 0 ; 10 ], fechado em ambos os extremos.
Para os gráficos de distribuição das notas das ¹⁵ Uma definição dessas estatísticas pode ser encontrada no Glossário. questões discursivas, foram consideradas mais duas categorias : questão em branco16 e nota zero. 3.1 CONDIÇÃO DE PRESENÇA E ESTATÍSTICAS BÁSICAS DA PROVA Esta seção apresenta a condição de presença, além das estatísticas selecionadas e histogramas da nota geral ( 3.1.1 ) e de cada componente : Formação Geral ( 3.1.2 ) e Conhecimento Específico ( 3.1.3 ).
São também apresentadas estatísticas selecionadas de subpopulações caracterizadas por Grande Região, Categoria Administrativa e Organização Acadêmica. 3.1.1 Condição de Presença e Estatísticas Básicas Gerais A Tabela 3.1 apresenta a Condição de Presença na prova, por Grande Região, dos estudantes concluintes de Administração.
A população total de Inscritos foi de 151.469.
Destes, 122.962 estiveram Presentes, sendo 18,8 \% o índice de não comparecimento.
A Região de maior abstenção foi a Centro-Oeste ( 28,1 \%, ) e a de menor abstenção foi a Sul ( 14,8 \% ).
Já a Tabela 3.2 apresenta as Estatísticas Básicas da prova, por Grande Região, dos estudantes concluintes de Administração.
A Média das notas da prova como um todo ( nas seções seguintes serão analisados os Componentes de Formação Geral e de Conhecimento Específico ) foi 41,7, sendo que os alunos da região Centro-Oeste obtiveram a Média mais baixa ( 38,8 ), e os da região Sudeste obtiveram a Média mais alta ( 42,5 ).
As demais médias foram : 40,9 na região Norte, 42,0 na região Nordeste e 41,3 na região Sul.
O Desvio padrão para o Brasil como um todo foi 12,2, sendo o maior Desvio padrão encontrado na região Nordeste ( 12,5 ) e o menor, na região Norte ( 11,0 ), indicando uma dispersão um pouco menor das notas desta última região. 16 Nesse grupo estão incluídas também as respostas classificadas como nulas ou desconsideradas.
A região que obteve a maior nota Máxima foi a Sudeste ( 90,9 ), ao passo que a região que atingiu a menor nota Máxima foi a Norte ( 78,5 ).
A Mediana do Brasil como um todo foi 40,8, sendo a maior Mediana obtida na região Sudeste ( 41,7 ), e a menor obtida no Centro-Oeste ( 37,9 ).
A nota Mínima foi zero em todas as regiões.
Considerando -se as notas segundo Grande Região, observa -se que existe diferença estatisticamente significativa ao nível de 95 \% entre a menor Média, obtida na região Centro-Oeste ( 38,8 ), e as médias de todas as demais regiões.
O comportamento das notas dos estudantes de todo o Brasil pode ser observado no Gráfico 3.1 que apresenta um histograma com a distribuição das mesmas.
Essa é uma distribuição unimodal com a moda no intervalo ( 30 ; 40 ].
A Tabela 3.3 apresenta informações referentes à Condição de Presença desagregadas por Categoria Administrativa e por Organização Acadêmica.
Da população total de Inscritos, 134.266 são de IES Privadas e 17.203 de IES Públicas.
Em relação à Organização Acadêmica, a maior participação foi obtida por estudantes de Universidades ( 54.578 ), seguido de Faculdades ( 50.218 ) e Centros Universitários ( 18.166 ).
A Categoria Administrativa de maior abstenção foi a Privada ( 18,8 \% ), e entre as Organizações Acadêmicas foi a dos Centros Universitários ( 21,3 \% ), que ficou acima da taxa de abstenção nacional de 18,8 \%.
Já a Tabela 3.4 apresenta informações referentes à Média da nota final desagregadas por Categoria Administrativa e por Organização Acadêmica.
A Média das notas da prova como um todo foi 41,7.
Em relação à Categoria Administrativa, os estudantes das IES Públicas obtiveram Média mais alta ( 48,3 ), e os das IES Privadas obtiveram Média mais baixa ( 40,9 ), que a Média nacional.
Observa -se que existe diferença estatisticamente significativa entre as médias das notas das IES Públicas e Privadas.
A diferença entre as médias das regiões Sudeste e Centro-Oeste ( 3,8 ), a maior e a menor Média, é inferior à diferença entre IES Públicas e Privadas ( 7,4 ), caracterizando uma maior diversidade administrativa do que regional.
No tocante à Organização Acadêmica, apenas as Universidades obtiveram Média mais alta ( 42,5 ) que a nacional.
As demais médias foram : 41,4 nos Centros Universitários e 41,0 nas Faculdades.
O Desvio padrão para as IES Públicas ( 13,6 ) e para as Universidades ( 12,8 ) foi superior ao do Brasil como um todo ( 12,2 ), indicando uma dispersão um pouco maior das notas nesta Categoria Administrativa e nesta Organização Acadêmica.
Constata -se que existe diferença estatisticamente significativa ao nível de 95 \% nas médias das notas dos estudantes dos três tipos de Organização Acadêmica. 3.1.2 Estatísticas Básicas no Componente de Formação Geral A Tabela 3.5 apresenta as Estatísticas Básicas em relação ao componente da prova que avalia a Formação Geral dos estudantes concluintes.
Os alunos de todo Brasil obtiveram desempenho médio de 53,3.
Quanto à variabilidade, o Desvio padrão das notas dos estudantes do Brasil, como um todo, foi 15,1.
A maior Média foi obtida na região Sudeste ( 53,8 ), e a menor, na região Centro-Oeste ( 50,9 ).
As demais médias foram : 53,3 na região Norte, 53,4 na região Nordeste e 53,0 na região Sul.
Já o maior Desvio padrão foi obtido na região Nordeste ( 15,7 ), e o menor, na região Norte ( 14,7 ).
Os demais desvios padrões foram : 15,0 na região Sudeste, 14,9 da região Sul, e 15,6 na região Centro-Oeste.
A maior nota no Componente de Formação Geral da prova do Enade/2015 foi 98,8, obtida por, pelo menos, um aluno na região Sudeste.
A menor nota Máxima foi obtida na região Centro-Oeste ( 97,4 ).
A Mediana do Brasil, como um todo, foi 53,8, sendo a menor Mediana encontrada na região Centro-Oeste ( 51,4 ) e a maior encontrada na região Sudeste ( 54,4 ).
A nota Mínima nesta parte foi zero em todas as regiões, sem exceção.
Considerando -se as notas segundo Grande Região, observa -se que existe diferença estatisticamente significativa entre a maior Média das notas do Componente de Formação Geral, obtida na região Sudeste ( 53,8 ), e na região Centro-Oeste ( 50,9 ), a menor Média.
O Gráfico 3.2 propicia a avaliação do desempenho dos estudantes do Componente de Formação Geral, a partir do histograma da distribuição das notas correspondentes.
A distribuição é unimodal, com moda em ( 50 ; 60 ], enquanto na prova, como um todo ( Gráfico 3.1 ), a moda foi alcançada dois intervalos abaixo ( 30 ; 40 ].
Nota -se, ainda que, no Gráfico 3.2, as notas apresentam maior dispersão do que no Gráfico 3.1 ( distribuição das notas da prova ), confirmado pela comparação dos desvios padrões : 12,2 para a nota da prova como um todo e 15,1 para o Componente de Formação Geral.
Na Tabela 3.6 são apresentadas as informações referentes ao desempenho dos concluintes do Componente de Formação Geral, em diferentes agregações : Categoria Administrativa e Organização Acadêmica.
Observa -se que existe diferença estatisticamente significativa entre as médias dos tipos de Categoria Administrativa.
A maior Média foi obtida por estudantes de IES Públicas ( 58,7 ), com uma diferença estatisticamente significativa da obtida por estudantes de IES Privadas ( 52,6 ).
Considerando -se o tipo de Organização Acadêmica, nota -se que há diferença estatisticamente significativa entre as médias dos tipos de Organização Acadêmica. 3.1.3 Estatísticas Básicas do Componente de Conhecimento Específico A Tabela 3.7 apresenta as Estatísticas Básicas referentes ao Componente de Conhecimento Específico da área de Administração.
A Média do desempenho dos alunos do Brasil, como um todo, foi 37,9.
A maior Média foi obtida na região Sudeste ( 38,8 ), e a menor, na região Centro-Oeste ( 34,7 ).
As demais médias foram : 36,7 na região Norte, 38,2 na região Nordeste e 37,3 na região Sul.
Quanto à variabilidade das notas, o Desvio padrão do Brasil, como um todo, foi 13,3, sendo o maior Desvio padrão observado na região Sudeste ( 13,6 ), e o menor, na região Norte ( 12,0 ).
Os demais desvios foram : 13,5 na região Nordeste, 13,1 na região Sul e 12,6 na região Centro-Oeste.
A Mediana das notas dos estudantes de todo o Brasil foi 36,7.
A maior Mediana ocorreu na região Sudeste ( 37,5 ), e a menor, na região Centro-Oeste ( 33,6 ).
As demais medianas foram : 35,9 na região Norte, 37,4 na região Nordeste e 36,1 na região Sul.
A nota Máxima do Brasil, como um todo, foi 95,1, sendo obtida por, pelo menos, um aluno na região Sudeste.
As demais notas máximas foram : 80,2 na região Norte, 92,7 na região Nordeste, 91,7 na região Sul e 83,7 na região Centro-Oeste.
A nota Mínima foi zero em todas as regiões.
Observa -se que existe diferença estatisticamente significativa da menor Média das notas do Componente de Conhecimento Específico da região Centro-Oeste ( 34,7 ) em relação às demais regiões.
O intervalo de confiança da região Norte ( 0,7 ) é significativamente maior do que de todas as outras regiões.
Assim como os Gráficos 3.1 e 3.2, o Gráfico 3.3, apresentado a seguir, permite uma avaliação do desempenho de concluintes de Administração em relação ao Componente de Conhecimento Específico com um histograma da distribuição das notas correspondentes.
Dentre as três distribuições apresentadas, esta é a mais concentrada nas notas baixas.
Esta também é uma distribuição unimodal, e o grupo modal é o ( 30 ; 40 ], o mesmo grupo modal da prova como um todo e abaixo do grupo modal para a Formação Geral.
A Tabela 3.8 apresenta uma comparação dos resultados em relação à Categoria Administrativa e à Organização Acadêmica, agora levando em conta o desempenho de alunos do Componente de Conhecimento Específico da prova.
No que se refere à Organização Acadêmica, a maior Média foi das Universidades ( 38,6 ), vindo a seguir a dos Centros Universitários ( 37,5 ) e, depois, a das Faculdades ( 37,1 ).
O maior Desvio padrão, e acima do valor para o Brasil, como um todo, foi o das Universidades ( 14,0 ).
As Universidades também obtiveram as maiores notas Máxima ( 95,1 ) e Mediana ( 37,4 ).
Os Centros Universitários obtiveram nota Máxima igual a 90,5 e as Faculdades obtiveram nota Máxima igual a 93,2.
As medianas foram : 36,5 nos Centros Universitários e 36,1 nas Faculdades.
A nota Mínima foi zero para todas as Organizações Acadêmicas.
Observa -se que existem diferenças estatisticamente significativas ao nível de 95 \% no Componente de Conhecimento Específico entre as notas das Universidades em relação aos Centros Universitários e Faculdades e também entre estas duas categorias.
Quanto à Categoria Administrativa, observa -se um comportamento semelhante àquele da parte de Formação Geral e à prova como um todo, ou seja, existe diferença estatisticamente significativa entre as médias das IES Públicas ( 44,8 ) e IES Privadas ( 37,0 ).
Neste caso também, a maior Média foi obtida por alunos de IES Públicas de ensino. 3.2 ANÁLISE DAS QUESTÕES OBJETIVAS Esta seção apresenta estatísticas selecionadas e histogramas das Questões Objetivas de Formação Geral ( 3.2.1 ).
São também apresentadas e comparadas as médias das subpopulações caracterizadas por Grande Região. 3.2.1 Componente de Formação Geral A Tabela 3.9 apresenta as Estatísticas Básicas relativas às oito questões objetivas do componente da prova que abrange a Formação Geral dos estudantes.
A Média do Brasil foi 50,5.
A menor Média foi encontrada na região Centro-Oeste ( 48,3 ), e a maior, na região Sudeste ( 50,9 ).
As demais médias foram : 49,4 na região Norte, 50,5 na região Nordeste e 50,7 na região Sul.
O Desvio padrão do Brasil foi 19,8, sendo o maior Desvio padrão encontrado na região Nordeste ( 20,2 ), e o menor, na região Sul ( 19,5 ).
Os demais desvios foram : 19,9 na região Norte, 19,8 na região Sudeste e 20,0 na região Centro-Oeste.
As medianas ( 50,0 ), as notas máximas ( 100,0 ) e as notas mínimas ( 0,0 ) foram iguais para todas as regiões.
A Tabela 3.10 apresenta o Índice de Facilidade e o Índice de Discriminação ( ponto bisserial ) para cada uma das questões objetivas do Componente de Formação Geral.
Quanto ao Índice de Facilidade, foram usadas as seguintes cores para diferenciar o nível de dificuldade da questão :  Azul para as questões classificadas com índice Muito fácil ( \&gt ; =0,86 ), verde para as questões classificadas com índice Fácil ( 0,61 a 0,85 ), amarelo para as questões classificadas com Médio ( 0,41 a 0,60 ), vermelho para as questões classificadas com Difícil ( 0,16 a 0,40 ) e roxo para as questões classificadas com Muito difícil ( \&lt ; =0,15 ).
Já quanto ao Índice de Discriminação, foram usadas as seguintes cores para qualificar a questão :  As questões classificadas com índice Fraco receberam a cor vermelho ( \&lt ; =0,19 ), as classificadas com Médio receberam a cor amarelo ( 0,20 a 0,29 ), as classificadas com Bom receberam a cor verde ( 0,30 a 0,39 ) e as classificadas com Muito bom ( \&gt ; =0,40 ) receberam a cor azul.
As questões objetivas do Componente de Formação Geral, segundo o Índice de Facilidade, foram assim avaliadas : das oito questões, nenhuma teve o Índice de Facilidade classificado como Muito fácil.
Três questões foram tidas como Fácil, por terem índice de acertos situado na faixa entre 0,61 e 0,85 ( de 61,0 \% a 85,0 \% de acertos ).
Duas questões foram consideradas de dificuldade Médio, situando -se no intervalo entre 0,41 e 0,60 do Índice de Facilidade, ou seja, houve entre 41,0 \% e 60,0 \% de acertos, enquanto três questões foram classificadas na categoria Difícil, situando -se no intervalo entre 0,16 e 0,40.
Por fim, nenhuma das questões apresentou menos de 15 \% de acertos, razão pela qual seria classificada como Muito difícil.
Como já comentado, para análise das questões objetivas, relativas à Formação Geral, segundo o poder de discriminação, utilizou -se o Índice de Discriminação ( ponto A Tabela 3.10 apresenta o Índice de Facilidade e o Índice de Discriminação ( ponto bisserial ) para cada uma das questões objetivas do Componente de Formação Geral.
Quanto ao Índice de Facilidade, foram usadas as seguintes cores para diferenciar o nível de dificuldade da questão :  Azul para as questões classificadas com índice Muito fácil ( \&gt ; =0,86 ), verde para as questões classificadas com índice Fácil ( 0,61 a 0,85 ), amarelo para as questões classificadas com Médio ( 0,41 a 0,60 ), vermelho para as questões classificadas com Difícil ( 0,16 a 0,40 ) e roxo para as questões classificadas com Muito difícil ( \&lt ; =0,15 ).
Já quanto ao Índice de Discriminação, foram usadas as seguintes cores para qualificar a questão :  As questões classificadas com índice Fraco receberam a cor vermelho ( \&lt ; =0,19 ), as classificadas com Médio receberam a cor amarelo ( 0,20 a 0,29 ), as classificadas com Bom receberam a cor verde ( 0,30 a 0,39 ) e as classificadas com Muito bom ( \&gt ; =0,40 ) receberam a cor azul.
As questões objetivas do Componente de Formação Geral, segundo o Índice de Facilidade, foram assim avaliadas : das oito questões, nenhuma teve o Índice de Facilidade classificado como Muito fácil.
Três questões foram tidas como Fácil, por terem índice de acertos situado na faixa entre 0,61 e 0,85 ( de 61,0 \% a 85,0 \% de acertos ).
Duas questões foram consideradas de dificuldade Médio, situando -se no intervalo entre 0,41 e 0,60 do Índice de Facilidade, ou seja, houve entre 41,0 \% e 60,0 \% de acertos, enquanto três questões foram classificadas na categoria Difícil, situando -se no intervalo entre 0,16 e 0,40.
Por fim, nenhuma das questões apresentou menos de 15 \% de acertos, razão pela qual seria classificada como Muito difícil.
Como já comentado, para análise das questões objetivas, relativas à Formação Geral, segundo o poder de discriminação, utilizou -se o Índice de Discriminação ( ponto bisserial ).
Nesta análise, as questões foram assim avaliadas : seis das oito questões apresentaram índice acima ou igual a 0,40 e, assim, foram classificadas com índice Muito bom para esse grupo de estudantes.
Duas questões tiveram Índice de Discriminação Bom, com valor entre 0,30 e 0,39, para esse grupo de estudantes.
Nenhuma questão teve nível Médio ou Fraco de discriminação para esse grupo de estudantes.
O Índice de Facilidade variou de 0,31 a 0,74, e o de Discriminação, de 0,34 a 0,48.
As seis questões com Índice de Discriminação Muito bom figuraram entre os diversos níveis de dificuldade desse conjunto : três classificadas na categoria Fácil ( questões 3, 5, e 6 ) do Índice de Facilidade, duas, na categoria Médio ( questões 2 e 8 ) e uma, na categoria Difícil ( questão 7 ).
Em particular, a questão 2 foi a que apresentou o maior poder discriminatório, com índice 0,48, porém foi considerada Médio em termos de facilidade, com uma proporção de 0,41 acertos.
O máximo de acertos foi alcançado pela questão 5 com um Índice de Facilidade de 0,74.
A questão de número 8 apresentou um Índice de Facilidade de 0,52, ou seja, mais da metade dos estudantes conseguiu resolvê -la, dentro do universo de participantes.
Seu Índice de Discriminação foi Muito bom ( 0,42 ).
Já a questão 7 obteve Índice de Discriminação Muito bom, 0,45, porém seu Índice de Facilidade foi Difícil ( 0,31 ). bisserial ).
Nesta análise, as questões foram assim avaliadas : seis das oito questões apresentaram índice acima ou igual a 0,40 e, assim, foram classificadas com índice Muito bom para esse grupo de estudantes.
Duas questões tiveram Índice de Discriminação Bom, com valor entre 0,30 e 0,39, para esse grupo de estudantes.
Nenhuma questão teve nível Médio ou Fraco de discriminação para esse grupo de estudantes.
O Índice de Facilidade variou de 0,31 a 0,74, e o de Discriminação, de 0,34 a 0,48.
As seis questões com Índice de Discriminação Muito bom figuraram entre os diversos níveis de dificuldade desse conjunto : três classificadas na categoria Fácil ( questões 3, 5, e 6 ) do Índice de Facilidade, duas, na categoria Médio ( questões 2 e 8 ) e uma, na categoria Difícil ( questão 7 ).
Em particular, a questão 2 foi a que apresentou o maior poder discriminatório, com índice 0,48, porém foi considerada Médio em termos de facilidade, com uma proporção de 0,41 acertos.
O máximo de acertos foi alcançado pela questão 5 com um Índice de Facilidade de 0,74.
A questão de número 8 apresentou um Índice de Facilidade de 0,52, ou seja, mais da metade dos estudantes conseguiu resolvê -la, dentro do universo de participantes.
Seu Índice de Discriminação foi Muito bom ( 0,42 ).
Já a questão 7 obteve Índice de Discriminação Muito bom, 0,45, porém seu Índice de Facilidade foi Difícil ( 0,31 ).
O Gráfico 3.4, para exemplificar, analisa o comportamento da questão de número 2 de Formação Geral.
Trata -se de uma questão considerada mediana, com relação à facilidade e a que obteve o maior Índice de Discriminação dessa parte da prova.
Neste gráfico, cada uma das cinco curvas representa o percentual de respostas em determinada alternativa da questão, em função do número de acertos dos estudantes nessa parte da prova ( Formação Geral/Múltipla Escolha ), antes de possíveis eliminações pelo critério do ponto bisserial.
A curva em preto corresponde à alternativa C, a correta para esta questão.
Assim, observa -se que entre os estudantes com menor número de acertos, nessa parte do exame, a situação mais frequente foi a escolha da alternativa D ( em roxo ), incorreta.
Por exemplo, entre os estudantes que acertaram três questões, 37,6 \% escolheram a alternativa D, 21,7 \% a alternativa C ( correta ), 15,4 \% a alternativa A, 14,7 \% a alternativa B e 10,1 \% a alternativa E.
Entre os que acertaram três respostas entre as questões de múltipla escolha de formação geral, 0,2 \% deixou a questão em branco e 0,3 \% marcou mais de uma alternativa, invalidando a questão.
À medida em que o número de acertos aumenta, indicando desempenho melhor nessa parte da prova, aumenta concomitantemente a proporção de estudantes que selecionaram a alternativa correta C, atingindo 100 \% para os estudantes com 8 acertos.
Essa análise permite verificar como a questão discriminou os grupos de desempenho, justificando o alto índice obtido na questão.
Cumpre notar que não é possível inferir deste gráfico nem o índice de facilidade ( que seria uma média da proporção ponderada pela quantidade de alunos com cada uma das notas ), nem o índice de discriminação ponto bisserial, por razão equivalente.
No caso extremo no qual a grande concentração dos acertos dos alunos fosse abaixo de 4, o índice de facilidade seria obrigatoriamente abaixo de 40 \% ( neste exemplo ).
Caso a concentração fosse em 6 acertos ou mais, o índice seria obrigatoriamente acima de 60 \%.
Os gráficos relativos às demais questões de Formação Geral constam do Anexo I. 3.2.2 Componente de Conhecimento Específico A Tabela 3.11 apresenta as Estatísticas Básicas em relação às questões objetivas do Componente de Conhecimento Específico da prova de Administração por Grande Região.
A Média do Brasil deste componente foi de 37,5.
A menor Média foi observada na região Centro-Oeste ( 34,6 ), e a maior, na região Sudeste ( 38,3 ).
O Desvio padrão de todo o Brasil foi 14,1, sendo o menor Desvio padrão encontrado na região Norte ( 12,9 ), e o maior, nas regiões Nordeste e Sudeste ( 14,3 ).
A Mediana de todo o Brasil foi 36,0, a mesma encontrada em todas as regiões, exceto na região Centro-Oeste que apresentou valor para a Mediana igual a 32,0.
A nota Máxima da prova foi 96,0, obtida nas questões objetivas do Componente de Conhecimento Específico, por, pelo menos, um aluno nas regiões Nordeste, Sudeste e Sul.
Nas demais regiões, a nota Máxima da prova foi 92,0 na região Norte e 88,0 na região Centro-Oeste.
Em todas as regiões a nota Mínima foi zero.
A Tabela 3.12 apresenta os Índices de Facilidade e Discriminação ( ponto bisserial ) das questões objetivas do Componente de Conhecimento Específico para os estudantes de Administração.
Para facilitar a diferenciação das questões, usaremos as mesmas cores da Tabela 3.10 para as diferentes classificações dos Índices de Facilidade e de Discriminação.
A partir do Índice de Facilidade obtido, pode -se concluir que um pouco mais de dois terços das questões objetivas da prova foi considerado, pelo menos, Difícil : das 27 questões válidas, 20 foram classificadas como Difícil ou como Muito difícil.
Uma questão foi classificada como Fácil, e outras seis consideradas como Médio.
Já quanto aos índices de discriminação das questões objetivas do Componente de Conhecimento Específico da prova, tem -se como resultado a seguinte classificação : 10 das 27 questões foram consideradas como boas, enquanto três delas tiveram Índice de Discriminação Muito bom.
Assim, para quase metade das questões – 13 em 27 – os Índices de Discriminação foram Bom ou Muito bom.
Dentre as demais, 12 delas foram classificadas como Médio e outras duas como Fraco, sendo 14, por conseguinte, a quantidade de questões nos dois patamares mais baixos de discriminação.
Constata -se, assim, que a prova – no que se refere ao Componente de Conhecimento Específico – possuía boa capacidade de discriminar entre aqueles que dominam ou não o conteúdo.
Dentre as questões que alcançaram os maiores índices de discriminação, três delas, as de números 20, 24 e 27, foram classificadas com Índice de Discriminação Muito bom, situando -se no intervalo de 0,40 a 0,50 do índice, todas foram classificadas na categoria Médio, quanto ao Índice de Facilidade.
A questão de número 25 foi a mais difícil dentre as 27 questões específicas, com baixo Índice de Facilidade, apenas 9,0 \% de acertos.
Essa questão apresentou poder discriminatório igualmente baixo, 0,06, o que comprova ter sido esta questão a mais difícil para os estudantes.
Destaca -se, também, a questão 19, com Índice de Facilidade 0,28, o que, em termos percentuais, corresponde a 28,0 \% de estudantes que responderam acertadamente.
Já 0,10 foi o seu de Índice de Discriminação.
Tais questões foram, portanto, pelo critério ponto bisserial, consideradas inadequadas.
Por isso, as questões 19 e 25 foram eliminadas do cômputo da nota final.
A título de exemplo das análises do comportamento das questões objetivas, o Gráfico 3.5 analisa a questão 24 do Componente de Conhecimento Específico.
Esta questão foi considerada pelas respostas dos estudantes avaliados na prova como uma questão com nível de facilidade Médio, 0,47, ou seja, 47,0 \% dos estudantes assinalaram acertadamente a opção B, correspondente ao gabarito.
Seu Índice de Discriminação foi igual a 0,44, classificado como Muito bom, também sendo esta questão a que apresentou o maior valor.
Neste gráfico, cada uma das cinco curvas representa o percentual de respostas em determinada alternativa da questão 24, em função do número de acertos dos estudantes nessa parte da prova, antes de possíveis eliminações de questões pelo critério do ponto bisserial.
A alternativa correta B, representada no gráfico pela curva em verde, foi escolhida em maiores proporções pelos alunos com desempenho melhor nessa parte da prova.
Já as alternativas incorretas, também denominadas distratores, foram selecionadas, principalmente, por aqueles com notas mais baixas.
Neste caso também a soma não é sempre 100 \% por causa das questões não respondidas ou com mais de uma opção marcada.
Aqueles com nota zero, na sua quase totalidade deixaram esta questão em branco ou marcaram mais de uma alternativa, comportamento considerado inválido.
A proporção de alunos que selecionou a resposta correta B aumenta gradativamente, chegando a atingir 100 \% para 24 acertos ou mais, enquanto a proporção dos que escolheram alternativas incorretas decai, a partir de dois acertos, como função do número de acertos nessa parte da prova.
Os gráficos relativos às demais questões do Conhecimento Específico constam do Anexo I. 3.3 ANÁLISE DAS QUESTÕES DISCURSIVAS Esta seção apresenta estatísticas selecionadas e histogramas das Questões Discursivas de Formação Geral ( 3.3.1 ) e Conhecimento Específico ( 3.3.2 ).
São também apresentadas e comparadas as médias de subpopulações caracterizadas por Grande Região. 3.3.1 Componente de Formação Geral As análises dos resultados de desempenho dos estudantes de Administração nas duas questões discursivas relativas à Formação Geral encontram -se na Tabela 3.13 e no Gráfico 3.6.
Na Tabela 3.13, observa -se que a nota Média nesse conjunto de questões foi acima da obtida nas objetivas.
Os estudantes, de todo o Brasil, obtiveram, em Formação Geral, Média 50,5 nas questões objetivas e 57,4 nas questões discursivas.
Pode -se notar também que o Desvio padrão nesse conjunto de questões foi próximo do obtido nas objetivas, 19,8 nas questões objetivas e 19,7 nas questões discursivas.
A maior Média foi obtida na região Norte ( 59,2 ), e a menor, na região Centro-Oeste ( 54,8 ).
A Mediana de todo o Brasil, neste componente, foi 62,0.
Nas regiões Norte, Nordeste e Sudeste a Mediana foi maior, 63,0, 63,0 e 62,5, respectivamente.
Nas demais regiões foi menor que o valor para o Brasil : Sul ( 61,0 ), Centro-Oeste ( 60,0 ).
A nota Máxima ( 96,0 ) foi a mesma nas regiões Nordeste e Centro-Oeste, sendo 99,0 na região Nordeste, 98,5 na região Sudeste e 97,0 na região Sul.
A nota Mínima ( 0,0 ) foi a mesma em todas as regiões do Brasil.
O Gráfico 3.6 representa a distribuição das notas nas questões discursivas do Componente de Formação Geral.
A moda desta distribuição ocorre no intervalo ( 60 ; 70 ].
Destaca -se também o intervalo [ 0 ; 10 ] com distribuição próxima a 6 \% do total de notas, sendo que no intervalo [ 0 ; 10 ] incluem -se além da nota zero, a frequência de alunos que deixaram este tipo de questão em branco.
Na sequência, os resultados verificados para cada uma das questões discursivas de Formação Geral serão apresentados, estabelecendo -se relações com os conteúdos abordados em cada uma delas.
Os comentários da Banca de docentes corretores, a respeito do observado na correção das respostas dos estudantes, suas impressões e conclusões serão apresentados junto à análise de cada questão.
Cumpre esclarecer que, tendo em vista que as questões discursivas de Formação Geral são padronizadas, ou seja, constam de todas as provas, os comentários da Banca são os mesmos para todas as carreiras acadêmicas, sendo direcionados a todos os estudantes que participaram do Enade/2015.
A seguir, serão analisados os desempenhos da Área de Administração nas duas questões discursivas de Formação Geral do Enade/2015, comparando -se os resultados obtidos com comentários para cada questão. 3.3.1.1 Análise de Conteúdo da Questão Discursiva 1 do Componente de Formação Geral Os dados de Administração, obtidos a partir das respostas à questão 1, encontram -se na Tabela 3.14 e no Gráfico 3.7.
Nessa questão – de melhor desempenho dentre as duas de Formação Geral – os alunos, de todo o Brasil, tiveram Média 62,3.
A maior Média para a questão 1 foi obtida na região Sudeste ( 63,1 ), e a menor, na região Centro-Oeste ( 59,6 ).
Quanto à variabilidade das notas, o Desvio padrão de todo o Brasil foi 21,3.
O menor Desvio padrão foi obtido na região Norte ( 20,3 ), e o maior Desvio padrão foi obtido na região Centro-Oeste ( 22,9 ).
As medianas de todas as regiões e do Brasil como um todo, foram iguais ( 65,0 ).
As notas máximas e mínimas da questão discursiva 1 também foram as mesmas para todas as regiões do Brasil, respectivamente, 100,0 e 0,0.
O Gráfico 3.7 mostra a distribuição das notas na questão discursiva 1 do Componente de Formação Geral.
Observa -se que a maior frequência corresponde aos alunos que obtiveram nota dentro do intervalo ( 70 ; 80 ], correspondendo à moda da distribuição.
Destacam -se, também, os alunos que deixaram a questão em branco com aproximadamente 7 \% do total. 3.3.1.2 Comentários sobre a correção de Conteúdo das respostas à Questão Discursiva 1 O enunciado da Discursiva 1 de Formação Geral estava claro e indicava um tema de extrema relevância, qual seja : a busca pela educação universal, especificamente, que apontava um caso onde ocorreu a discriminação de acesso ao ensino formal atingindo uma menina, representando a categoria das mulheres discriminadas em tantos campos da vida e, mais ainda, em sociedades que cultuam o poder soberano dos homens.
A linguagem utilizada foi clara, indicando a existência dos fatos vinculados à luta da menina Malala pelo acesso ao ensino formal das jovens do seu país, e o confronto ideológico que gerou o ataque à sua integridade física.
Diante da repercussão internacional, houve o reconhecimento pelo prêmio Nobel, e a ONU resolveu apoiá -la como símbolo da luta pela igualdade entre os sexos buscando a universalização da educação.
O conteúdo está integrado às Diretrizes, vez que o acesso ao ensino formal está vinculado ao direito constitucional à educação, que, também, é protegido pelo Estatuto da Criança e do Adolescente e nos pactos de Direitos Humanos aderidos pelo Brasil.
Por outro lado, o texto demandava uma reflexão sobre a realidade dos indivíduos que não têm acesso à educação e a situação das mulheres excluídas desse processo por segmentos da sociedade.
Como se tratava de texto contendo uma ampla gama de possibilidades, era natural a ocorrência de interpretações diversas das esperadas.
Releva notar que o enunciado pretendeu, primordialmente, identificar a questão de gênero, apontando para as dificuldades que as mulheres possuem, não somente no Brasil, mas em vários países do mundo.
Ao analisar as provas realizadas, constatamos que no item ‘ a ’ poucos foram os que identificaram os estatutos normativos indicados no padrão de resposta.
Essa ausência pode ser justificada tendo em vista que, no enunciado indicado, não existia elemento que pudesse ser utilizado pelo concluinte para embasar sua resposta nesse caminho.
As respostas foram mais gerais, identificando que a educação constitui um item fundamental para o desenvolvimento do indivíduo e que as mulheres sofrem mais com as limitações impostas por países com regime patriarcal, que discriminam o sexo feminino e, com frequência, proíbem as mulheres de frequentar a escola.
Deve -se ter em vista que a indicação de resposta que vinculava o texto aos estatutos normativos representou um raciocínio próprio dos estudantes de Direito que, naturalmente, encaminharam suas respostas para a fundamentação nessa linha identificando, não somente a Constituição, mas também leis, regulamentos, etc.
Tal linha de argumentação não foi objeto da formação dos estudantes das demais áreas.
A maior parte dos concluintes estabeleceu o tema educação formal associado à busca pelo conhecimento, como acesso ao ensino, indicando a independência intelectual e o desenvolvimento da consciência crítica, limitando as manipulações e permitindo amplo acesso ao mercado de trabalho, traduzindo uma libertação econômica e ascensão social.
No tocante ao item ‘ b ’, restou claro que a maioria esmagadora das respostas confluiu para as indicações do padrão apresentado, sendo relevante o tema de igualdade de gênero, além de candente para as concluintes do sexo feminino que desenvolveram os textos de forma concatenada, fundamentada e, muitas vezes, lamentando que, nos dias atuais, ainda se discuta a necessidade de igualar os gêneros na sociedade.
Nessa linha, foram identificadas respostas, consoante com o padrão apresentado, especialmente nos seguintes temas : • violência física e psicológica contra a mulher, incluindo a Lei Maria da Penha, no caso específico do Brasil ; • tolerância/intolerância a vestimentas, trajes, comportamentos socialmente estereotipados ; • aspectos socioculturais que impõem à mulher uma condição de submissão na sociedade, tais como a impossibilidade de a mulher manifestar seus desejos e posicionamentos em algumas culturas, entre outros ; • igualdade/desigualdade de gênero, por exemplo no mercado de trabalho, em relação à desigualdade salarial ; • ideais de Liberdade, Igualdade e \textbf{Fraternidade} : referência a esses ideais como possibilidade de equilibrar as relações de poder entre homens e mulheres.
Em muitos casos, também, os concluintes fizeram um liame entre o acesso à educação e a melhora das relações entre homens e mulheres ao permitir que a conclusão de maior grau de educação formal permita o acesso a cargos públicos relevantes, empregos bem remunerados e posições de poder, tanto na área pública quanto na área privada.
Como as opções do segundo item foram mais amplas, o padrão de resposta identificou as linhas de pensamento adotadas pelos concluintes.
A maior parte das respostas identificou, claramente, o problema da relação de gênero e suas consequências na sociedade.
Também foi frequente constatarem que existe efetiva melhora na situação das mulheres no Brasil, mas que tal movimento não ocorre em todo o mundo.
Aliás, o exemplo de Malala conduziu a uma reflexão em âmbito mundial, e os concluintes apresentaram suas análises críticas à situação da mulher em países que não admitem sequer a sua inserção no plano educacional e, muito menos, no mercado de trabalho.
A maior parte dos concluintes apresentou concatenação lógica nas suas respostas e uma minoria, pouco expressiva, não demonstrou esse domínio de 3.3.1.3 Análise de Conteúdo da Questão Discursiva 2 do Componente de Formação Geral A Tabela 3.15 mostra que o desempenho médio dos estudantes na questão discursiva 2 ( média 55,0 ) foi inferior ao obtido na questão discursiva 1 ( média 62,3 ).
A região Norte foi aquela cuja Média, nessa questão, foi maior ( 56,2 ), e a de menor Média foi a região Centro-Oeste ( 52,3 ).
Quanto à variabilidade das notas, o Desvio padrão de todo o Brasil foi 23,0, superior ao obtido na questão discursiva 1 ( 21,3 ).
O maior desvio nessa questão foi obtido na região Nordeste ( 24,3 ), enquanto o menor foi obtido nas regiões Sudeste e Sul ( 22,6 cada ).
A Mediana de todo o Brasil foi 61,0, a mesma das regiões Norte, Nordeste, Sudeste.
A Mediana encontrada na região Sul foi 60,0 e na região Centro-Oeste foi 59,0.
As notas máximas ( 100,0 ) foram as mesmas em todas as regiões do Brasil, exceto na região Norte, que obteve nota Máxima 99,0 e região Centro-Oeste, com nota Máxima 98,0.
As notas mínimas ( 0,0 ) também foram as mesmas em todas as regiões, sem exceção.
O Gráfico 3.8 mostra a distribuição das notas na questão discursiva 2 do Componente de Formação Geral.
Observa -se que a maior frequência corresponde aos alunos que obtiveram nota dentro do intervalo ( 60,70 ].
Destacam -se, também, os alunos que deixaram a questão em branco e os alunos que obtiveram nota dentro do intervalor ( 20,30 ], caracterizando -se como máximos locais.
Nota -se, ainda, que as notas ficaram mais distribuídas ao longo dos intervalos em comparação à questão discursiva de número 1.
Isto pode ser constatado também pela comparação do desvio padrão das notas da questão discursiva 2 ( 23,0 ) e da questão discursiva 1 ( 21,3 ). 3.3.1.4 Comentários sobre a correção de Conteúdo das respostas à Questão Discursiva 2 O enunciado da Discursiva 2 estava, igualmente, claro, indicando um tema de extrema relevância, qual seja : o reconhecimento de manifestações artísticas nascidas na periferia, como integrante da cultura reconhecida pela sociedade.
A linguagem utilizada foi clara, indicando a análise dos fatos e suas consequências no cotidiano dos partícipes e o confronto com outras visões culturais que são consideradas mais adequadas pela sociedade civil.
O conteúdo está integrado às Diretrizes, vez que o aspecto cultural e sua inter-relação com o conhecimento haurido no nível superior é essencial para a compreensão de mundo do universitário concluinte, com reflexo na sua atuação profissional.
O texto demanda uma reflexão sobre a realidade dos indivíduos que não têm acesso aos itens básicos da sociedade, como educação, saúde, cultura tradicional e segurança e que, apesar disso, criam alternativas para o seu lazer.
A incorporação do funk como manifestação cultural e artística da sociedade brasileira encontrou, basicamente, duas correntes de resposta : • os que aquiescem positivamente à integração ; • os que rejeitam, acerbamente, tal possibilidade.
No meio, uma corrente que somente admite a integração com modificações na conduta dos integrantes do movimento, com a exclusão de letras exclusivamente com apelos sexuais e incitação à violência.
A interpretação positiva indica a origem da música na periferia, notadamente nas comunidades cariocas e sua projeção para o Brasil e o mundo.
Nessa perspectiva, surge o reconhecimento como cultura popular comparada ao samba, axé e outros ritmos também nascidos no Brasil.
Também ocorre comparação com o movimento hip-hop americano.
Já a análise negativa prioriza a vinculação do ritmo com traficantes, prostitutas, vilipêndio às mulheres, falta de controle quanto ao som, local dos eventos, transações sexuais em público, uso frequente de drogas, violência \textbf{exacerbada}.
Faltam, para essa corrente, elementos básicos da cultura, tratando -se de música sem autonomia como movimento de mudança ou aspecto próprio.
A corrente intermediária reconhece o vínculo originário com as comunidades dominadas por quadrilhas de traficantes, mas admite o seu reconhecimento de acordo com a mudança dos atores que transitam na modalidade, como músicos, \textbf{intérpretes}, para adequar as letras das canções expurgando os aspectos considerados \textbf{nocivos} como o apelo sexual, as indicações positivas para a violência e apologia às drogas.
Um aspecto bastante criticado nas músicas dessa modalidade foi a exploração da mulher, entendida como um objeto de satisfação masculina.
É relevante tal constatação, pois a questão pertinente à Malala reflete a luta pela valorização da mulher.
Os concluintes compreenderam a ideia que foi problematizada e escolheram uma estratégia adequada para solucioná -la.
O caráter dinâmico da cultura foi comumente expressado, não através do enfrentamento abordado no texto, mas mediante a indicação ( de forma rasa e genérica ) das dificuldades e consequente aceitação social de manifestações culturais, antes marginalizadas e que hoje integram o acervo cultural nacional, tais como : samba, grafite, capoeira e forró.
Quanto às reflexões de Laraia, a grande maioria reconheceu o funk como manifestação cultural, porém, o identificou mais na sua expressão musical, no balanço de seu ritmo, tecendo contundentes críticas, e mesmo forte discordância quanto ao conteúdo e mensagem das letras quando remetem à depreciação da imagem feminina ou induzem à apologia ao crime.
De resto, os textos abordaram bem mais o item dois do padrão de Respostas.
Uma minoria enfrentou os temas mudança e dinamismo cultural.
As respostas da corrente positiva indicam que o preconceito sempre acompanhou as manifestações populares como o samba, na sua origem.
Admitem que diante da miscigenação surgem movimentos culturais de várias naturezas, o que inclui o funk.
Também indicam a dificuldade de lazer dos membros das comunidades o que faz surgirem eventos para preencher essa necessidade.
Nesse âmbito, apontam a música como alternativa de ascensão social e integração aos bens de consumo.
As respostas da corrente negativa vinculam a música ao tráfico de drogas cujos elementos comandam as comunidades, propiciando que os bailes sejam mecanismos de venda e consumo dessa mercadoria ilícita.
Na mesma toada, indicam a vinculação com a exploração sexual, principalmente das mulheres maiores e menores de idade.
Indicam aspectos imorais nas letras.
Apresentam os locais dos bailes como inadequados e provocadores da \textbf{desordem} pública.
Recusam que músicas com discriminação às mulheres sejam consideradas como cultura.
A corrente intermediária indica a possibilidade de adoção do funk como cultura desde que ocorram modificações nas letras e no comportamento dos atores ( músicos, letristas e público ), retirando os aspectos negativos já apontados.
Fala -se em “ funk do bem ”.
A maior parte das respostas apresentadas indica que os concluintes têm conhecimento dos conflitos sociais e, no caso da questão do funk, conseguem identificar a origem popular da música.
A partir daí é que ocorre a cisão entre os que defendem o movimento como cultura, aduzindo a ascensão social dos artistas envolvidos dos que vinculam a música com atos ilícitos, rejeitando sua inserção no âmbito cultural.
Em algumas provas, há dificuldade de concatenação lógica de raciocínio, mas a maior parte consegue desenvolver o texto de forma razoável.
Essa questão permite visualizar a divisão social com os pensamentos \textbf{liberais} e conservadores, sendo bastante claros nas respostas dos concluintes. 3.3.1.5 Análise de Língua Portuguesa das Questões Discursivas do Componente de Formação Geral Os dados de Administração, obtidos a partir das respostas às questões discursivas do Componente de Formação Geral, no que tange à Língua Portuguesa, encontram -se na Tabela 3.16 e no Gráfico 3.9.
Nesse aspecto, os alunos, de todo o Brasil, tiveram Média 56,9.
A maior Média com respeito à Língua Portuguesa foi obtida na região Norte ( 58,8 ), e a menor, na região Centro-Oeste ( 54,3 ).
Quanto à variabilidade das notas, o Desvio padrão de todo o Brasil foi 20,3.
O menor Desvio padrão foi obtido na região Norte ( 19,3 ) e o maior Desvio padrão foi obtido na região Centro-Oeste ( 21,5 ).
A Mediana das notas de Língua Portuguesa foi 61,0 para o Brasil como um todo.
A maior Mediana com respeito à Língua Portuguesa foi obtida na região Norte ( 62,5 ), e a menor, na região Centro-Oeste ( 59,5 ).
A nota Máxima para todo o Brasil foi de 99,5, com, pelo menos, um aluno tirando essa nota na região Nordeste.
Além disso, a nota Mínima foi zero em todas as regiões do país.
O Gráfico 3.9 mostra a distribuição das notas de Língua Portuguesa do Componente de Formação Geral.
Observa -se que a maior frequência ( quase 30 \% ) corresponde aos alunos que obtiveram nota no intervalo ( 60 ; 70 ].
Destaca -se, também, os alunos que deixaram ambas as questões em branco representando pouco mais de 5 \% do total, caracterizando -se como máximo local. 3.3.1.6 Comentários sobre a correção das respostas de Formação Geral com respeito à Língua Portuguesa Ao encaminhar as questões 1 e 2 na direção da produção de um texto dissertativo, espera -se que o participante utilize seus conhecimentos sobre o assunto e estruture seus textos de acordo com as exigências do registro formal próprio dessa situação comunicativa.
Essa configuração determina exigências quanto : à adequação da seleção vocabular, ao desenvolvimento do conteúdo, à estruturação sintática dos períodos, à organização lógica das ideias, à utilização de procedimentos de encadeamento textual e referenciação, à obediência às exigências morfossintáticas próprias da modalidade escrita da norma-padrão, ao respeito às regras ortográficas e às regras de acentuação gráfica.
O padrão de resposta utilizado na avaliação das questões 1 e 2 considerou os aspectos relevantes ao bom desempenho linguístico como competências distintas, de modo a permitir um mapeamento detalhado do domínio dos recursos disponíveis na Língua Portuguesa para a comunicação escrita formal.
Com base nesse objetivo, foram avaliados os seguintes aspectos : a ) Estruturação textual condizente com o gênero solicitado ( texto dissertativo ) e o modo de organização textual expositivo adequado ao gênero – essa competência envolve :  distribuição do conteúdo do texto em parágrafos, de modo a garantir a sua organização temática ;  estruturação sintática condizente com o padrão da modalidade escrita formal da língua portuguesa de modo a garantir a clareza necessária ;  utilização de operadores discursivos que contribuam para a progressão temática do texto, estabelecendo relações lógicas entre as ideias apresentadas, tanto do ponto de vista intrafrasal, como do interfrasal ;  utilização de procedimentos de referenciação lexical e pronominal que permitam a retomada de referentes textuais ;  utilização de sinais de pontuação que contribuam para a organização lógica da frase e do texto.
Espera -se, portanto, que o participante recorra a procedimentos linguístico-discursivos para organizar seu texto, permitindo o encadeamento lógico entre suas partes de forma a garantir a progressão e a coerência textual.
Isso significa que os seguintes procedimentos foram penalizados, de acordo com o padrão de resposta proposto :  elaboração de frases fragmentadas que comprometam a estrutura lógico-gramatical do texto ;  sequência justaposta de ideias sem encaixamentos sintáticos, reproduzindo hábitos da oralidade ;  elaboração de frase com apenas oração subordinada, sem oração principal ;  emprego equivocado do conector ( preposição, conjunção, pronome relativo, alguns advérbios e locuções adverbiais ) que não expresse a relação lógica adequada entre dois trechos do texto e prejudique a compreensão da mensagem ;  repetição ou substituição inadequada de palavras sem se valer dos recursos oferecidos pela língua ( pronome, advérbio, artigo, sinônimo ) ;  utilização inadequada dos sinais de pontuação comprometendo a clareza textual. b ) Respeito às convenções ortográficas da norma-padrão da Língua Portuguesa – essa competência envolve o domínio das regras de acentuação gráfica e da grafia padrão das palavras ( com ausência de abreviaturas próprias da linguagem da internet ), de acordo com as convenções estabelecidas pela legislação em vigor e consubstanciadas no Vocabulário Ortográfico da Língua Portuguesa, editado pela Academia Brasileira de Letras ( com aceitação da legislação anterior, no caso das regras relativas ao uso do hífen e da acentuação gráfica, já que o exame foi realizado no mês de \textbf{novembro} de 2015 ).
A competência relativa ao respeito às convenções ortográficas envolve :  grafia correta das palavras ;  respeito às regras de acentuação gráfica ;  emprego de maiúsculas em início de frase, em nomes próprios de pessoas, lugares ou instituições ;  ausência de abreviações como p/, vc, tb, pq, tá, né, usadas muitas vezes em escrita informal e na internet ;  obediência às regras de separação de sílabas no final da linha. c ) Domínio dos diferentes aspectos morfossintáticos próprios da modalidade escrita formal da norma-padrão da Língua Portuguesa – os itens relacionados a essa competência são : a concordância nominal, a concordância verbal, a regência nominal, a regência verbal, a flexão nominal, a flexão verbal, a correlação entre os tempos verbais, a colocação pronominal.
O domínio dessa competência envolve :  flexão do verbo para estabelecer concordância de número com o sujeito da frase ;  flexão do artigo, do adjetivo e do pronome para concordar em número e em gênero com o substantivo a que se referem ;  respeito à regência nominal e à verbal, utilizando a preposição adequada depois de um substantivo, de um verbo ou de um adjetivo ; emprego da preposição antecedendo o pronome relativo e emprego adequadamente os pronomes relativos “ cujo(a) ” e “ onde ” ;  emprego adequado do acento grave indicador da crase entre uma preposição e um artigo ( a+a ) ;  atendimento às regras de colocação pronominal ( próclise e ênclise ), distintas dos hábitos da oralidade ou da escrita informal. d ) Seleção vocabular adequada à modalidade escrita formal da Língua Portuguesa, exigida pela situação comunicativa – essa competência envolve a precisão na utilização do vocabulário relacionado à temática solicitada pela questão ; a ausência de marcas da oralidade, como termos de sentido muito genérico ( “ coisa ”, “ negócio ”, “ você ” ) e termos de registros mais informais ( como gírias, jargões, frases feitas, ditados populares, termos regionais ).
Assim, espera -se que o participante respeite a adequação vocabular não empregando gírias ou expressões coloquiais, evite repetição desnecessária de palavras e utilize um vocabulário mais formal, como solicitado por um texto dissertativo.
A escolha dessas competências para subsidiar o processo de avaliação apoia -se na concepção de que, no desempenho dos graduandos, a modalidade escrita tem apresentado uma intensa simplificação, originada no padrão da modalidade oral da Língua Portuguesa.
No caso do texto de base dissertativa, inscrito em um registro formal, a distância entre as duas modalidades é ainda maior, o que provoca situações de hipercorreção ( desvios provocados pela utilização inadequada de uma regra da norma-padrão ) e de truncamentos sintáticos ( estruturas frasais incompreensíveis devido à complexidade sintática própria da modalidade escrita ).
Observam -se, então, os seguintes aspectos que marcam essa distinção entre as duas modalidades, devido à excessiva simplificação da modalidade falada : a ) redução drástica de estruturas subordinadas, compensada pelo aumento na frequência de estruturas coordenadas e absolutas, por um lado, ou pela elaboração de estruturas truncadas pelo excesso de ideias sem a devida conexão subordinativa ; b ) redução no uso de operadores argumentativas para expressar relações lógicas essenciais à construção do texto, substituídas pela exigência de inferência por parte do interlocutor para suprir a sua ausência ; c ) redução do uso do subjuntivo, ao lado da ampliação do uso do indicativo combinado a estruturas frasais coordenadas ou absolutas ; d ) empobrecimento do processo de referenciação, com a repetição de pronomes ou nomes ; e ) simplificação extrema da marcação da categoria tempo na morfologia verbal ; f ) falta de domínio de vocabulário mais abstrato e de maior complexidade, essencial ao desenvolvimento do processo dissertativo ; g ) redução drástica no emprego da acentuação gráfica, processo intensificado pela intensa utilização das redes sociais.
Os aspectos macroestruturais da elaboração do texto não foram avaliados neste processo porque dizem respeito à avaliação do conteúdo.
Portanto, a banca de Formação Geral, composta por profissionais de diferentes áreas do conhecimento, se encarregou da avaliação do atendimento ao solicitado no enunciado das questões do ponto de vista do desenvolvimento do conteúdo.
A avaliação do desempenho linguístico considerou, portanto, três grandes grupos de competências, segundo os aspectos explicitados anteriormente : a ) domínio das convenções ortográficas : grafia de vogais e consoantes, uso de maiúsculas e minúsculas, emprego do hífen e acentuação gráfica ; b ) domínio dos procedimentos de estruturação textual do ponto de vista microestrutural : organização interna dos períodos, emprego de conectores para a articulação lógica entre os períodos e entre os parágrafos, emprego de marcas de referenciação lexical e pronominal ; utilização dos sinais de pontuação que contribuem para a organização lógica da frase. c ) domínio das regras de caráter morfossintático estabelecidas como modelares do ponto de vista da modalidade escrita formal da norma-padrão da Língua Portuguesa : concordância nominal e verbal, regência nominal e verbal, colocação pronominal, flexão nominal e verbal, correlação entre tempos e modos verbais, ausência de marcas de oralidade.
A seleção vocabular adequada à modalidade escrita formal da Língua Portuguesa foi incorporada à essa última competência, tendo em vista a intersecção entre as duas do ponto de vista das exigências do registro formal da modalidade escrita da norma-padrão.
Apresenta -se, a seguir, o padrão de resposta aprovado pelo Inep e já utilizado na avaliação do desempenho linguístico das questões 1 e 2 de Formação Geral no Enade/2013 e no Enade/2014.
Às competências, reunidas nos três grupos descritos anteriormente, foram atribuídos os seguintes \textbf{pesos} relativos : aspectos ortográficos ( 20 \% ) ; aspectos textuais ( 40 \% ), aspectos morfossintáticos e vocabulares ( 40 \% ).
Com base na avaliação do processo desenvolvido nos dois anos anteriores, utilizaram -se quatro níveis para a avaliação do desempenho linguístico nas três competências.
O nível zero ficou reservado para casos especiais, como fuga total ao tema, palavras soltas, produção de uma frase incompleta, entre outros.
A avaliação das respostas elaboradas pelos participantes revelou um resultado coerente na comparação entre as duas questões.
Entretanto, em relação aos resultados do Enade/2014, observou -se uma melhora do domínio de todas as competências, talvez devido ao perfil dos cursos analisados, como Direito e Jornalismo, que pressupõem a utilização profissional da linguagem.
Observou -se uma diferença no desempenho dos participantes em relação às duas questões discursivas, provavelmente em função dos temas solicitados : - a Questão 1 solicitava que o participante desenvolvesse um texto dissertativo sobre o significado da premiação de Malala, abordando dois aspectos : o direito das jovens à educação formal e as relações de poder entre homens e mulheres no mundo.
Esse encaminhamento favoreceu o desenvolvimento mais articulado do texto, o que se refletiu em melhor desempenho dos aspectos vocabulares e textuais. - a Questão 2 solicitava que o participante desenvolvesse um texto dissertativo para discutir a questão do reconhecimento do funk como legítima manifestação artística e cultural do povo brasileiro.
Por ser um tema muito polêmico, o tema do preconceito sócio-cultural-econômico favoreceu a construção de textos mais críticos, dividindo o conjunto de participantes em dois grupos : os que defendiam esse tipo de manifestação artística como legítima e os que a criticavam de maneira veemente, chegando, inclusive, a utilizar muitos termos de baixo calão.
Essa última atitude teve como consequência a elaboração de textos mal desenvolvidos do ponto de vista sintático-discursivo, com muitos truncamentos.
Os resultados revelaram essa diferença.
Ao analisar os dados de cada questão, constata -se que a Questão 1 teve maior número de questões de médio ou bom desempenho, em relação à Questão 2.
Por outro lado, a Questão 2 apresentou um índice maior de notas baixas em relação à Questão 1.
Os resultados da avaliação correspondem aos seguintes aspectos observados em cada competência : a ) aspectos ortográficos : o desempenho dos participantes revelou uma diferença muito grande nos dois aspectos analisados : baixo índice de desvios ortográficos e grande índice de desvios de acentuação.
Em vários casos, ocorre ausência completa de acentuação gráfica.
Os resultados revelam que a tendência dominante entre os universitários brasileiros é a eliminação da acentuação gráfica, talvez motivada pelos hábitos relacionados às redes sociais e pela ausência de esclarecimento dos meios de comunicação, das autoridades e das escolas sobre as decisões do Acordo Ortográfico da Língua Portuguesa de 1990.
Os casos mais sistemáticos de eliminação do acento indicador da sílaba tônica são :  palavras proparoxítonas ( “ proximos ”, “ politicos ” ) ;  palavras paroxítonas terminadas em ditongo crescente ( “ necessario ”, “ noticiarios ”, “ individuo ”, “ dependencia ”, “ varios ”, “ propria ” ) ;  palavras oxítonas ( “ ninguem ”, “ esta ”, “ ate ”, “ tambem ”, “ alvara ” ) ;  uso indevido do acento gráfico em palavras como “ genêro ”, “ melâncolia ”, por exemplo.
Quanto ao domínio das convenções relativas à grafia das palavras, observam -se desvios como : I. hipercorreção pela escolha de “ e ” no lugar de “ i ” ou de “ i ” no lugar de “ e ”, por influência de hábitos da oralidade ( “ descriminação ” por “ discriminação ”, “ entervenção ” por “ intervenção ”, “ esteriotipada ” por “ estereotipada ” ) ; II. eliminação do “ r ” marcador do infinitivo verbal ( “ esta ” no lugar de “ estar ”, “ estuda ” no lugar de “ estudar ” ).
III. eliminação de uma sílaba na palavra “ educação ”, que foi grafada “ edução ” por um grande número de participantes.
IV. outros casos de desvios de grafia relacionados à variação diastrática podem ser observados em “ estrupo ”, “ subjulgadas ”, “ precoseito ”, “ soubre ”, “ vecendo ”, “ indesencia ”, “ apolojia ”, “ fulga ”.
V. desvios de segmentação podem ser atribuídos à confusão entre o \textbf{morfema} gramatical “ mos ” e o pronome oblíquo “ nos ”, em “ chegar mos ”, “ enfrentar -mos ”.
VI. inadequação no uso da maiúscula para destacar determinadas palavras-chave do texto, como “ Violência ”, “ Brasileiros ”.
Destaque -se, também, o grande número de participantes que grafam os textos inteiramente em caixa alta.
Vale observar, também, que, ao contrário do que se esperava, não apareceram abreviaturas próprias do “ internetês ”, relacionado ao uso de redes sociais e emails. b ) aspectos textuais : esta competência é a que se revela como a mais problemática entre os participantes, porque são muitos os problemas observados, desvios acumulados durante toda a formação do estudante : sequência justaposta de ideias sem encaixamentos sintáticos ; redução drástica de estruturas subordinadas, ao lado do aumento na frequência de estruturas coordenadas e absolutas ; redução no uso de conectores para expressar relações lógicas essenciais à construção do texto, substituídas pela exigência de inferência por  palavras paroxítonas terminadas em ditongo crescente ( “ necessario ”, “ noticiarios ”, “ individuo ”, “ dependencia ”, “ varios ”, “ propria ” ) ;  palavras oxítonas ( “ ninguem ”, “ esta ”, “ ate ”, “ tambem ”, “ alvara ” ) ;  uso indevido do acento gráfico em palavras como “ genêro ”, “ melâncolia ”, por exemplo.
Quanto ao domínio das convenções relativas à grafia das palavras, observam -se desvios como : I. hipercorreção pela escolha de “ e ” no lugar de “ i ” ou de “ i ” no lugar de “ e ”, por influência de hábitos da oralidade ( “ descriminação ” por “ discriminação ”, “ entervenção ” por “ intervenção ”, “ estereotipada ” por “ estereotipada ” ) ; II. eliminação do “ r ” marcador do infinitivo verbal ( “ esta ” no lugar de “ estar ”, “ estuda ” no lugar de “ estudar ” ).
III. eliminação de uma sílaba na palavra “ educação ”, que foi grafada “ edução ” por um grande número de participantes.
IV. outros casos de desvios de grafia relacionados à variação diastrática podem ser observados em “ estrupo ”, “ subjulgadas ”, “ precoseito ”, “ soubre ”, “ vecendo ”, “ indesencia ”, “ apolojia ”, “ fulga ”.
V. desvios de segmentação podem ser atribuídos à confusão entre o \textbf{morfema} gramatical “ mos ” e o pronome oblíquo “ nos ”, em “ chegar mos ”, “ enfrentar -mos ”.
VI. inadequação no uso da maiúscula para destacar determinadas palavras-chave do texto, como “ Violência ”, “ Brasileiros ”.
Destaque -se, também, o grande número de participantes que grafam os textos inteiramente em caixa alta.
Vale observar, também, que, ao contrário do que se esperava, não apareceram abreviaturas próprias do “ internetês ”, relacionado ao uso de redes sociais e emails. b ) aspectos textuais : esta competência é a que se revela como a mais problemática entre os participantes, porque são muitos os problemas observados, desvios acumulados durante toda a formação do estudante : sequência justaposta de ideias sem encaixamentos sintáticos ; redução drástica de estruturas subordinadas, ao lado do aumento na frequência de estruturas coordenadas e absolutas ; redução no uso de conectores para expressar relações lógicas essenciais à construção do texto, substituídas pela exigência de inferência por parte do interlocutor para suprir a sua ausência ; emprego equivocado de operadores que não estabelecem relações lógicas coerentes entre ideias do texto ; emprego inadequado do pronome relativo ( com omissão da preposição ou a utilização de pronome inadequado, como “ onde ” ) ; repetição exaustiva de termos sem a utilização de procedimentos mais sofisticados de substituição ( hiperonímias, hiponímias, nominalizações, expressões metafóricas ) ; frases fragmentadas que comprometem a estrutura lógico-gramatical ; frases formadas apenas por oração subordinada, sem oração principal. parte do interlocutor para suprir a sua ausência ; emprego equivocado de operadores que não estabelecem relações lógicas coerentes entre ideias do texto ; emprego inadequado do pronome relativo ( com omissão da preposição ou a utilização de pronome inadequado, como “ onde ” ) ; repetição exaustiva de termos sem a utilização de procedimentos mais sofisticados de substituição ( hiperonímias, hiponímias, nominalizações, expressões metafóricas ) ; frases fragmentadas que comprometem a estrutura lógico-gramatical ; frases formadas apenas por oração subordinada, sem oração principal.
Um importante aspecto a destacar é o baixíssimo desempenho de uma parte dos participantes em relação à estrutura formal do texto produzido, o que é extremamente preocupante ao se levar em conta que são graduandos em fase final de formação.
São frequentes os casos de desvios de estruturação frasal, com uso inadequado ou ausência de conectivos entre parágrafos e entre frases.
Em uma parte dos textos, falta um mínimo de textualidade e de domínio do registro padrão da língua.
Na verdade, observam -se relações linguísticas quase agramaticais, como as estabelecidas pela sequência de gerúndios sem o apoio de um ponto de partida para a organização das informações gramaticais e semânticas.
Observou -se que uma grande parte dos participantes não organizou as ideias em parágrafos, talvez devido ao pequeno número de linhas disponibilizadas para a resposta da questão ou, quem sabe, pela suposição de que não seria necessária essa divisão por não se tratar de um texto no modelo de uma redação dissertativo-argumentativa, como solicitado nos vestibulares.
Quanto à utilização dos mecanismos de referenciação, deve -se destacar a ocorrência, em uma boa parte dos textos, de repetições de palavras ou expressões sem a utilização de termos sinônimos ou pronomes, como seria adequado.
Outro aspecto relevante a ser destacado é a quase total ausência de operadores argumentativos, tanto intrafrasais como interfrasais, repercutindo uma tendência atual da mídia escrita.
Quanto à utilização dos sinais de pontuação, observou -se uma grande precariedade nos textos analisados.
É muito frequente a ocorrência de parágrafos sem marca interna de pontuação para separar os períodos.
Vale observar que não foi penalizada a ausência de vírgula para destacar locuções ou adjuntos adverbiais de pequena extensão deslocados de posição na frase, por ser um uso opcional.
São os seguintes os tipos de problemas encontrados : I. vírgula : utilização de vírgula para separar o sujeito e o predicado ; ocorrência de apenas uma das vírgulas para separar uma palavra, uma expressão ou uma oração encaixada ; uso de vírgula no lugar do ponto para separar ideias que constituem períodos distintos ; II. ponto e vírgula : utilização do ponto e vírgula no lugar de vírgula ; III. ponto final : ausência de ponto final para separar períodos. uma oração encaixada ; uso de vírgula no lugar do ponto para separar ideias que constituem períodos distintos ; II. ponto e vírgula : utilização do ponto e vírgula no lugar de vírgula ; III. ponto final : ausência de ponto final para separar períodos. c ) aspectos morfossintáticos e vocabulares : os resultados são muito transparentes em relação aos aspectos mais problemáticos no desempenho dos participantes.
O desvio mais frequente, em relação à regência, é a falta do sinal indicativo da crase – isso revela que o usuário não tem consciência de que, sob a forma do termo “ a ”, existe a presença de uma preposição “ a ”, exigida pela regência do termo anterior, e de um artigo definido.
Embora em outros exames, como o Enem, a falta de crase seja penalizada em acentuação, nesta avaliação esse desvio foi considerado no âmbito dos aspectos morfossintáticos.
Outro problema relacionado à regência verbal e à nominal, encontrado frequentemente nas questões, foi a ausência de preposição antes de pronome relativo, processo generalizado na modalidade oral da língua, em situações de registro informal.
Apesar da possibilidade de que essa alteração de regência se generalize no padrão escrito da Língua Portuguesa, como já está ocorrendo até em textos jornalísticos, o não emprego da preposição foi penalizado neste processo de avaliação.
Outro desvio muito frequente diz respeito aos processos de concordância verbal e de concordância nominal.
Quanto à concordância de número, observou -se ausência de marca com sujeito posposto ou uso indevido ( uso inadequado da marca de plural comandado pelo núcleo plural da locução adjetiva, apesar de o substantivo que funciona como núcleo do sintagma nominal estar no singular ).
Uma ocorrência generalizada foi a ausência de acento circunflexo na forma plural do presente do indicativo do verbo “ ter ”, que foi considerada como um desvio na concordância verbal e não na acentuação gráfica.
Quanto à concordância de gênero, vários casos foram observados, normalmente no âmbito de sintagmas nominais longos, em que o adjetivo está afastado do substantivo.
Quanto à questão da colocação pronominal, foram poucos os desvios observados.
Apesar de serem aspectos relacionados à oralidade, concluiu -se que, no registro escrito formal, a maioria dos participantes já incorporou regras como a não introdução da frase por um pronome oblíquo e a próclise na presença de um termo atrator.
Não se adotou, entretanto, o padrão excessivamente formal descrito pelas gramáticas normativas em relação à posição do pronome oblíquo em locuções verbais, já que esse uso está muito distante da prática cotidiana, até em textos mais formais.
Quanto aos aspectos vocabulares, alguns tipos de inadequação foram observados, principalmente na Questão 2, relacionada ao funk : expressões da oralidade, seleção vocabular incompatível com o contexto, gerando situações de falta de inteligibilidade ; falta de domínio de vocabulário mais abstrato e de maior complexidade, essencial ao desenvolvimento do texto de base dissertativa.
Ainda em função da temática desenvolvida na Questão 2, foi significativa a presença de palavras de baixa calão, por retratar um universo cultural específico.
Os textos que continham esse tipo de vocabulário foram considerados inadequados em função do registro formal exigido e foram avaliados com grau zero.
Várias marcas de oralidade foram identificadas, embora não em alta frequência : o uso do pronome relativo “ onde ” como relativo universal, falta de artigo definido antes de substantivo, repetição de palavras por falta de vocabulário, reduções como “ tá ” e “ pra ”, expressões informais. 3.3.2 Componente de Conhecimento Específico Na parte da prova relativa às questões discursivas do Componente de Conhecimento Específico ( Tabela 3.17 ), observa -se que a Média foi bem mais baixa do que para as questões discursivas do Componente de Formação Geral.
Enquanto no Componente de Formação Geral a Média para estudantes de Administração de todo o Brasil foi 57,4, na parte de Conhecimento Específico a Média foi 40,0.
A maior Média deste componente foi obtida pelos estudantes da região Sudeste ( 41,5 ), e a menor, pelos da região Centro-Oeste ( 35,4 ).
Quanto à variabilidade das notas, o Desvio padrão de todo o Brasil foi 23,1.
O maior Desvio padrão foi encontrado nas regiões Nordeste e Centro-Oeste ( 23,2, cada ), e o menor, na região Norte ( 22,3 ).
A maior nota Máxima, 100, foi obtida em quase todas as regiões do país, exceto na região Centro-Oeste que obteve nota Máxima igual a 96,7.
Além disso, a nota Mínima ( 0,0 ) foi obtida por pelo menos um aluno em todas as regiões do Brasil.
A Mediana do Brasil como um todo foi 41,7.
A maior Mediana foi obtida na região Sudeste ( 43,3 ), e a menor, na região Centro-Oeste ( 36,7 ).
O Gráfico 3.10 representa a distribuição das notas nas questões discursivas no Componente de Conhecimento Específico.
A moda desta distribuição ocorre no intervalo ( 40 ; 50 ].
Destaca -se, também, os alunos que obtiveram nota no intervalo [ 0,10 ] com aproximadamente 13 \% do total.
A análise de cada uma destas questões será feita a seguir. 3.3.2.1 Análise da Questão Discursiva 3 do Componente de Conhecimento Específico Na questão 3, cujos resultados aferidos encontram -se descritos na Tabela 3.18, a Média dos estudantes de todo o Brasil foi 42,6.
A menor Média nessa questão foi obtida pelos alunos da região Centro-Oeste ( 38,9 ), enquanto a maior Média foi obtida na região Sudeste ( 43,8 ).
Quanto à variabilidade das notas, o Desvio padrão de todo o Brasil foi 27,4.
O maior Desvio padrão foi obtido na região Nordeste ( 28,1 ), enquanto o menor foi obtido na região Sul ( 26,7 ).
A nota Máxima, 100,0 pontos, foi alcançada por, pelo menos, um aluno de todas as regiões.
A Mediana do Brasil como um todo foi 40,0, a mesma para quase todas as regiões, exceto para a região Centro-Oeste que obteve a Mediana igual a 35,0.
A nota Mínima ( 0,0 ) foi obtida por, pelo menos, um estudante em todas as regiões do Brasil.
O Gráfico 3.11 mostra a distribuição das notas na questão discursiva 3, do Componente de Conhecimento Específico da área de Administração.
Sem levar em consideração que aproximadamente 13 \% dos estudantes deixaram a questão em branco, essa distribuição é unimodal com moda no intervalo ( 30,40 ]. 3.3.2.2 Comentários sobre as respostas à Questão Discursiva 3 O tema abordado na questão estava alinhado com o que é proposto nas Diretrizes Curriculares Nacionais para os cursos de graduação em Administração.
Além disso, a sustentabilidade é um tema bastante relevante para a sociedade.
Pelo padrão de resposta, o nível de profundidade esperado pareceu adequado a concluintes de cursos de Administração.
O enunciado apresentava um problema conceitual.
No entanto, o texto retirado do IBGE traz um erro a respeito da conceituação de consumerismo.
Esse conceito foi tratado de forma negativa e antagônica à sustentabilidade, o que é um equívoco crasso.
O consumerismo é definido na literatura de marketing como resultado de movimentos de consumidores focados em consumo consciente, racional e sustentável, que leva em conta as consequências dos atos de consumo.
O referido texto do IBGE deveria ter utilizado o termo consumismo, relativo a ações de consumo excessivo e que desprezam quaisquer consequências.
A despeito desse equívoco conceitual, as interpretações dos estudantes do que lhes era solicitado não levaram a respostas que destoassem do padrão de resposta adotado.
Ressalta -se que foram desconsiderados erros derivados do uso incorreto do termo consumerismo no enunciado.
A questão pareceu adequada para o tempo de prova, mas o espaço certamente foi reduzido para que fosse redigida uma resposta completa de todos os três itens solicitados.
Em que pese uma possível sobreposição entre os itens ‘ b ’ e ‘ c ’ – oportunidades de negócios para empresas podem estar relacionadas ao papel dessas na conservação do meio ambiente –, o espaço disponível para a resposta era limitado.
Dada a atualidade do tema e sua permanente discussão, não só no meio educacional, como também em vários veículos de mídia, a questão deveria ser considerada fácil.
No entanto, houve dificuldade na apresentação completa das dimensões da sustentabilidade, geralmente reduzida aos aspectos ambientais, assim como alguma dificuldade de argumentação completa a respeito do papel dos consumidores na conservação do meio ambiente.
Dessa forma, considerando a formação específica em Administração, a dificuldade foi média.
O padrão de respostas e a distribuição de pontos apresentaram orientações objetivas que permitiram a avaliação adequada das respostas dos estudantes.
Com relação ao item ‘ a ’, ficou claro o desconhecimento dos estudantes com relação às três dimensões da sustentabilidade.
O item comumente utilizado como base para essa discussão, o meio-ambiente, foi referido em praticamente todas as \textbf{repostas}.
No entanto, as dimensões social e econômica foram praticamente esquecidas.
Os itens ‘ b ’ e ‘ c ’ também não fugiram das discussões mais rasteiras e cotidianas a respeito de sustentabilidade.
Em muitos casos, os estudantes não conseguiam trazer respostas com o papel dos consumidores na conservação do meio ambiente.
Em outros, não havia propostas de oportunidades de negócios para as empresas e foi comum a repetição de trechos do enunciado como forma de preencher o espaço destinado à resposta.
A linguagem adotada pelos estudantes, via de regra, apresentou problemas da forma escrita da língua portuguesa ( ortografia e gramática ), o que também deixou às claras as deficiências de formação.
Além disso, havia dificuldade de encadeamento correto das ideias nas respostas.
Os estudantes apresentaram dificuldade de escrever um texto dissertativo que fosse assertivo e tivesse estrutura consistente.
Os textos mal escritos exigiam mais de uma leitura para serem compreendidos, nem sempre plenamente.
A dificuldade de expressão está clara nos textos fragmentados e com erros gramaticais ( principalmente de pontuação ) que dificultam a compreensão por parte do leitor.
Também fica clara a deficiência de formação do estudante, que não consegue se fazer entender plenamente.
Vale a pena apontar que a questão objetiva número 14 trazia certa sobreposição com o que deveria ser discutido nessa questão discursiva, ao tratar de engenharia reversa.
Alguns estudantes mais atentos citaram o processo de logística reversa como oportunidade de negócio para empresas.
Também parece interessante notar que alguns fatos relevantes, quando da aplicação da prova, como o desastre de Mariana e a crise hídrica em São Paulo, foram citados nas provas.
Isso pode mostrar que alguns estudantes demonstraram capacidade de associar parte do conteúdo aprendido no curso com a realidade que os cerca. 3.3.2.3 Análise da Questão Discursiva 4 do Componente de Conhecimento Específico A Tabela 3.19 contém as informações relativas à questão 4 do conjunto de questões do Componente de Conhecimento Específico.
O desempenho dos estudantes nessa questão foi inferior ao das questões 3 e 5.
A Média geral do Brasil foi 31,9, sendo a menor Média registrada na região Centro-Oeste ( 28,0 ), e a maior, na região Sudeste ( 33,3 ).
A nota Máxima ( 100,0 ) foi atingida por, pelo menos, um concluinte de todas as regiões.
A Mediana em todo o Brasil foi 30,0, o mesmo valor foi obtido em todas as regiões, exceto na região Centro-Oeste que obteve a Mediana igual a 15.
As notas mínimas foram zero em todas as regiões.
O Gráfico 3.12 representa a distribuição de notas da questão discursiva 4, do Componente de Conhecimento Específico.
Essa distribuição é multimodal com moda principal no intervalo ( 40 ; 50 ], seguida de perto pelo intervalo dos alunos que deixaram a questão em branco e pelo intervalo ( 10 ; 20 ], cada um representando mais de 17 \% do total. 3.3.2.4 Comentários sobre as respostas à Questão Discursiva 4 As Diretrizes Curriculares Nacionais para os cursos de graduação em Administração listam, dentre os conteúdos de formação profissional, em seu Art. 5o, a Administração de Recursos Humanos, aqui entendida como Gestão de Pessoas.
Esse conteúdo é fundamental para a formação do administrador, mesmo para aqueles que não tratarão diretamente da Gestão de Pessoas, uma vez que não existe função do administrador que possa ser realizada sem a interação entre pessoas.
O tema multiculturalidade na Gestão de Pessoas é pertinente e encontra eco em várias questões vivenciadas em organizações transnacionais na atualidade.
Diversas empresas brasileiras atuam no exterior, assim como diversas organizações estrangeiras atuam no Brasil.
Nesses ambientes de trabalho é comum que empregados advindos de diferentes realidades culturais tenham que conviver e desenvolver projetos em conjunto.
Os desafios para que isso aconteça de forma harmoniosa são muitos e incluem dificuldades de comunicação e de relacionamento, associadas a diferenças de valores, hábitos, crenças e comportamentos.
Por conta disso, os administradores devem acompanhar os momentos de interação entre funcionários para que os objetivos da organização possam ser alcançados.
A questão da globalização e da internacionalização de empresas foi um tema oportuno, mas grande parte dos estudantes não pareceu ter tido contato com essa realidade, conceitualmente ou como experiência vivida.
Em grande número, as respostas mostraram que os estudantes observaram alguma relação entre diversidade ( gênero, raça etc. ) com a realidade que percebiam ou viviam, no entanto, o conceito de multiculturalidade não foi desenvolvido de forma adequada em praticamente nenhuma resposta.
A percepção a respeito de uma atuação organizacional que \textbf{ultrapassa} limites nacionais ainda pareceu muito baixa.
O enunciado não se apresentou com objetividade necessária para que o estudante conseguisse ter uma ideia clara a respeito do que deveria responder.
Os trechos retirados de obras de Maximiniano e Cox, utilizados como introdução à questão, não ajudaram a entender o que estava sendo perguntado e foram pouco citados nas respostas.
Enquanto o texto de Maximiniano é uma opção bibliográfica frequente nas ementas dos cursos de graduação em Administração, o mesmo não pode ser dito a respeito do texto de Cox, em inglês.
Por conta disso, durante as provas do ENADE/2015, deve ter ocorrido o primeiro contato da grande maioria dos estudantes com a classificação das organizações proposta por esse último autor – monolíticas ; plurais ; multiculturais.
Dessa forma, pouquíssimos estudantes utilizaram esses conceitos em suas respostas ou souberam elaborar reflexões que envolvessem essa abordagem específica.
O padrão de respostas acordado com a Comissão Assessora levou em consideração as características do enunciado e foi ajustado para que as respostas dadas fossem avaliadas corretamente.
A partir da leitura das respostas apresentadas, ficou claro que os estudantes tiveram dificuldade na compreensão da questão.
Pareceu que a principal dificuldade se deveu ao fato da questão não se apoiar em uma objetividade mais latente ( não tratou de um caso que possa ser discutido mais objetivamente ).
A solicitação de se redigir um texto dissertativo também pareceu muito problemática por não auxiliar o aluno, conduzindo -o para responder ao que foi pedido, e por não deixar claro o nível de profundidade exigido.
Tendo em vista o espaço disponível e a amplitude da pergunta, as respostas foram muito vagas.
Por conta disso, considera -se que essa questão tenha sido difícil em comparação às demais questões discursivas de Conhecimentos Específicos.
Houve certa dificuldade por parte dos alunos em associarem o tema sugerido para a dissertação – benefícios e desafios de uma organização multicultural para a gestão de pessoas –, e os aspectos a serem elencados para responder à questão : a ) valorização da diversidade cultural e benefícios para a gestão de pessoas ; e b ) proposta de mudança organizacional que envolva diferentes áreas funcionais da organização.
Como os aspectos a serem tratados não lidavam com desafios da organização multicultural, e a pontuação da questão estava relacionada a esses itens, os estudantes tenderam a incluir em suas respostas exclusivamente os itens pontuados.
As respostas trouxeram mais discussões sobre a valorização da diversidade do que os desafios, e nem todos os estudantes chegaram a abordar o item ‘ b ’ em suas respostas.
Houve o uso de argumentos do senso comum e pouco suporte no conhecimento da Área para discutir a diversidade e o multiculturalismo nas organizações.
Além da deficiência teórica, os estudantes demonstraram deficiências no uso da língua escrita.
Apesar da solicitação de produção de textos dissertativos, de maneira geral, as respostas se apresentaram muito desestruturadas.
Os estudantes tiveram muita dificuldade em produzir um texto com conteúdo claro e estrutura consistente.
No tocante ao conteúdo mais específico da questão, alguns pontos deviam ser destacados.
Em primeiro lugar, ficou claro que a maioria dos estudantes não compreendeu o que seria a atuação global de uma organização e não considerou a expansão além dos limites locais.
Além disso, grande parte dos estudantes ignorou que a questão tratava de Gestão de Pessoas.
Muitos trouxeram as vantagens do multiculturalismo para outros setores da organização, como, por exemplo, o marketing, apontando a importância do funcionário que compreende o consumidor : “ gente diferente atende clientes melhor ” ; “ cada um com sua peculiaridade, passa a ser uma vantagem no atendimento ” ; “ empatia com consumidores ” ; “ melhora imagem no mercado ”.
De forma mais próxima da discussão de Gestão de Pessoas, alguns apontaram que o multiculturalismo está associado ao aumento da criatividade : “ gente diferente trabalhando gera inovação ” ; “ habilidades de um completar o outro ”.
No tocante ao item ‘ b ’, poucos trataram desse item em suas respostas.
Os que o fizeram, trouxeram a ideia de que uma ação interessante seria fazer com que o funcionário passasse por vários setores da empresa, ou trocasse sistematicamente de posição hierárquica para que cada membro compreendesse a posição do outro.
Além disso, em menor número, foram citadas as ações de recrutamento e seleção que levassem em conta a multiculturalidade.
Outras respostas foram muito genéricas e propuseram ações menos específicas.
As respostas consideradas fracas traziam : repetições de trechos do enunciado, chavões que não tinham consistência ; obviedades que pouco se relacionavam com o tema da questão e tratavam de dificuldades e benefícios de uma organização multicultural.
As respostas medianas foram as que conseguiram apresentar benefícios e desafios da multiculturalidade, mas não trouxeram propostas de mudança organizacional que envolvesse diversas áreas funcionais da organização.
Nas respostas consideradas boas os estudantes conseguiram contemplar todos os pontos esperados pelo padrão, o conceito de multiculturalidade era tratado corretamente, os benefícios e desafios da integração cultural eram abordados e era proposta alguma mudança organizacional consistente.
De maneira geral, as respostas apresentadas traziam uma introdução que tratava superficialmente do multiculturalismo, usando algum trecho repetido do enunciado, e, em seguida, um trecho no qual o estudante apresentava um benefício e um obstáculo para a organização multicultural.
Quase sempre os aspectos relacionados à Gestão de Pessoas eram deixados de lado e considerações genéricas a respeito de traços culturais, tais como alimentação ou hábitos religiosos, eram traçadas.
Poucas respostas traziam propostas de mudança organizacional que envolvesse as diferentes áreas funcionais da organização, o que denotou a dificuldade de o estudante relacionar o mundo à sua volta, problemas organizacionais e o conhecimento da Área de Administração. 3.3.2.5 Análise da Questão Discursiva 5 do Componente de Conhecimento Específico A Tabela 3.20 contém as informações relativas à questão 5 do conjunto do Componente de Conhecimento Específico.
O desempenho dos estudantes de todo o Brasil nesta questão foi mais alto dentre as três questões discursivas desse componente.
A nota Média dos estudantes de todo o Brasil foi 45,5.
A maior Média foi registrada na região Sudeste ( 47,5 ), enquanto a menor Média foi registrada na região Centro-Oeste ( 39,4 ).
Quanto à variabilidade das notas, o Desvio padrão dos alunos do Brasil, como um todo, foi 36,9.
Enquanto o maior desvio foi encontrado nas regiões Nordeste e Sudeste ( 36,9, cada ), o menor foi encontrado na região Norte ( 36,0 ).
A Mediana ( 35,0 ) e a nota Mínima ( 0,0 ) são as mesmas em todas as regiões.
Para o conjunto de alunos de Administração do Brasil, a nota Máxima foi 100,0, sendo obtida por, pelo menos, um estudante de cada uma das regiões.
O Gráfico 3.13 representa a distribuição das notas da questão discursiva 5 do Componente de Conhecimento Específico.
Essa distribuição é multimodal com moda principal no intervalo ( 60 ; 70 ], seguida de perto pelo intervalo ( 30 ; 40 ].
Destaca -se também, o grande número de estudantes que deixaram esta questão em branco ou que obtiveram nota no intervalo ( 90,100 ], cada um correspondendo a quase 20 \% do total. 3.3.2.6 Comentários sobre as respostas à Questão Discursiva 5 A questão apresentada era simples, com enunciado claro e comando objetivo, e, portanto, não gerou dúvidas de interpretação nos estudantes que realizaram a prova.
O nível de profundidade foi adequado para concluintes de cursos de Administração.
A questão estava adequada ao tempo de prova, e o espaço dedicado à resposta foi suficiente.
Dada sua objetividade e clareza, a questão pode ser considerada fácil quando comparada com as demais questões discursivas de Conhecimentos Específicos da prova de Administração.
Conforme as Diretrizes Curriculares Nacionais para os cursos de graduação em Administração, a análise de oportunidades e ameaças, presente na disciplina de Administração estratégica, encaixa -se como conteúdo necessário para a formação do administrador.
A referência utilizada na questão ( Barney e Hesterly, 2011 ) é uma obra relevante, que está presente em bibliografias de muitos cursos de graduação brasileiros e, portanto, é de conhecimento amplo.
Além disso, do ponto de vista do desenvolvimento de profissionais capacitados para atuação no mercado, as análises de oportunidades e ameaças são muito importantes para o direcionamento estratégico de organizações dos mais variados tipos.
O padrão de resposta estava bem definido, e as respostas dadas pelos estudantes, salvo quando completamente equivocadas, traziam até três oportunidades ou ameaças relacionadas com o ambiente da organização, dentre os previstos.
A avaliação foi bastante objetiva, o que se refletiu na distribuição de notas dessa questão.
A linguagem adotada pelos estudantes foi adequada.
Optaram por apresentar respostas em tópicos e não em forma de textos dissertativos.
Isso faz sentido frente às demandas do ambiente de trabalho, cada vez mais focado em comunicação rápida e simplificada, que exija pouco tempo para ser compreendida.
Dado o caráter mais objetivo da questão, os estudantes apresentaram desempenho mais satisfatório em comparação com as duas outras questões discursivas.
Pareceu haver maior facilidade de compreensão do comando nessa questão, assim como pareceu haver maior facilidade na apresentação de respostas, provavelmente por serem mais objetivas.
As respostas apresentadas traziam reflexões a respeito de possíveis oportunidades ou ameaças associadas aos elementos do ambiente geral listados no enunciado.
Parte dos estudantes fez reflexões mais abstratas com relação aos elementos do ambiente geral – “ tendência cultural é outro fator importante, pois hábitos e costumes dos compradores devem ser analisados... ” –, enquanto outros utilizaram situações específicas, como exemplos – “ A inovação é necessária, a Kodak, por exemplo, foi pioneira em \textbf{câmeras} fotográficas, mas perdeu o mercado diante dessa mudança tecnológica que a empresa não acompanhou. ” O uso de exemplos mostra a necessidade de uma referência objetiva por parte dos alunos, para que analisem melhor a situação que lhes é apresentada.
Em alguns casos, houve confusão na associação de elementos do ambiente.
Variáveis demográficas e geográficas foram confundidas, assim como houve alguma confusão envolvendo acontecimentos internacionais específicos e condições legais e políticas.
Nessa questão, o desempenho dos estudantes foi melhor.
A questão traz em seu enunciado a lista de elementos do ambiente, o que facilitou consideravelmente a circulação teórica dos estudantes.
Contudo, em alguns casos, ainda houve o uso de expressões corriqueiras na explicação do impacto desses elementos sobre as organizações ou a mera repetição do enunciado.
Os estudantes que têm dificuldade em escrever de forma dissertativa, utilizaram a estrutura de tópicos na resposta, o que foi possibilitado na medida em que o enunciado da questão pede a escolha de três 3.3.3 Considerações Finais As questões dissertativas que compuseram a parte de conhecimentos Específicos da prova de Administração cobriram três tópicos importantes para a formação do administrador : sustentabilidade, gestão de pessoas e gestão estratégica.
Dessa forma, o conteúdo proposto foi adequado para avaliar a qualidade da formação dos concluintes.
No entanto, uma questão apresentou problema em seu enunciado, que foi minimizado com a adaptação do padrão de resposta.
Conforme já mencionado nesse relatório, a Questão 3 trouxe imprecisão na conceituação de consumerismo, o que, provavelmente, trouxe alguma confusão para os estudantes.
As respostas a essa questão ficaram concentradas na dimensão ambiental da sustentabilidade, praticamente desconsiderando as dimensões social e econômica da sustentabilidade.
A Questão 4 trouxe alguns conceitos que podem não ser familiares a muitos dos estudantes de Administração ( organizações monolíticas, plurais e multiculturais ), e isso pode ter causado dificuldade para a formulação das respostas dos estudantes.
As respostas a essa questão demonstraram que os estudantes têm dificuldade em compreender a realidade global da Administração, assim como a deficiência conceitual a respeito de multiculturalidade e de diversidade no ambiente de trabalho.
A Questão 5 teve enunciado mais direto e, com isso, parece ter trazido mais conforto ao estudante.
Nesse caso, as respostas tenderam a ser diretas e, por vezes, sintetizadas por meio da estrutura de tópicos.
De forma geral, conforme já apontado, os estudantes tiveram muita dificuldade de expressão escrita, o que denotou deficiência no uso da norma padrão da Língua Portuguesa, que deve ser superada na formação do administrador.
A dificuldade de utilizar os conceitos das disciplinas também esteve clara.

% matched lemmas: câmera, desordem, exacerbar, fraternidade, intérprete, liberal, morfema, nocivo, novembro, peso, repor, ultrapassar
\end{textsample}
